%-----------------------------------------------------------------
% KATA PENGANTAR
%
% Sesuaikan daftar ucapan terima kasih dengan kondisi Anda. Nama, kota, dan
% tanggal pada blok tanda tangan diambil otomatis dari include/identitas.tex.
%-----------------------------------------------------------------
\preface % \preface JANGAN DIHAPUS!

Segala puji dan syukur kepada Tuhan Yang Maha Esa atas segala rahmat dan
karunia-Nya sehingga penulis dapat menyelesaikan penulisan Tesis yang berjudul
\textbf{``\MakeUppercase{\judul}''} sesuai dengan rencana. Selama penulisan
tesis ini penulis memperoleh berbagai arahan, bimbingan, dan bantuan dari
berbagai pihak. Oleh karena itu, melalui tulisan ini penulis mengucapkan rasa
terima kasih kepada:

\begin{enumerate}\setlength{\itemsep}{0pt}
    \item{Bapak/Ibu \dekan, Dekan \faculty{} \university, atas kesempatan dan dukungan yang diberikan.}
    \item{Bapak/Ibu \kajur, Ketua Departemen \dept, atas ilmu pengetahuan dan dukungan administratif selama perkuliahan.}
    \item{Bapak/Ibu \kaprodi, Ketua Program Studi \programstudi, atas segala fasilitas dan arahan selama studi.}
    \item{Bapak/Ibu \firstsupervisor, Dosen Pembimbing I, atas bimbingan, arahan, dan dukungan yang luar biasa selama penelitian hingga penulisan tesis ini selesai.}
    \item{Bapak/Ibu \secondsupervisor, Dosen Pembimbing II, atas masukan berharga dan dukungan dalam penyusunan tesis ini.}
    \item{Seluruh Dosen di Departemen \dept{} \faculty, atas ilmu pengetahuan, pengalaman, dan ketelitiannya selama perkuliahan.}
    \item{Orang tua dan keluarga tercinta, atas segala doa, dukungan moral maupun material, serta motivasi yang menjadi kekuatan terbesar penulis.}
    % Tambah daftar ucapan terima kasih lain di bawah baris ini
    \item{Teman-teman seperjuangan program studi \programstudi, atas kebersamaan, dukungan, dan semangat selama studi.}
\end{enumerate}

\vspace{0.2cm}

Penulis menyadari segala ketidaksempurnaan yang terdapat di dalamnya, baik dari
segi materi, cara, maupun bahasa yang disajikan. Seiring dengan itu, penulis
memharap kritik dan saran dari pembaca yang bersifat membangun demi
kesempurnaan tesis ini. Harapan penulis, semoga tesis ini dapat bermanfaat bagi
banyak pihak dan bagi perkembangan ilmu pengetahuan, khususnya bidang
\programstudi.

%-----------------------------------------------------------------
% Blok tanda tangan. Dipisah dengan \par agar tidak dapat menimpa
% paragraf penutup di atasnya, dan memakai \vspace positif supaya
% jaraknya tetap terjaga meski halamannya sudah hampir penuh.
%-----------------------------------------------------------------
\par\vspace{0.55cm}

% Catatan: memakai single spacing agar blok tanda tangan tetap rapi.
% Jarak setelah tanggal disediakan untuk tanda tangan basah.
\begin{singlespace}
\begin{tabular}{p{7.5cm}l}
	&\city, \disahkantanggal\\[1.5cm]
	&\underline{\fullname} \\
	&NIM. \idnum
\end{tabular}
\end{singlespace}
